\chapter{Formal Verifiable Provenance (v1.3.0)}
The ultimate defense against accusations of "fake results" or "overfitting" is cryptographic transparency. DTEAM v1.3.0 implements the "No-Fake" architecture.

\section{The Execution Manifest}
Every completed discovery epoch generates an \texttt{ExecutionManifest}, $M$:
\begin{equation}
M = \left\{ H(L), \pi(a_1, a_2, \dots, a_n), H(N) \right\}
\end{equation}
\begin{itemize}
    \item $H(L)$: The deterministic canonical hash (Merkle root equivalent) of the input \texttt{EventLog}.
    \item $\pi$: The exact sequence of $n$ RL actions taken by the agent.
    \item $H(N)$: The canonical topological hash of the output Petri net.
\end{itemize}

Because DTEAM is strictly deterministic, providing $M$ alongside a dataset allows any third party to re-run the engine in Reproduction Mode. If the engine executes $\pi$ against $L$, it is mathematically guaranteed to produce $H(N)$. The result is cryptographically anchored to the execution path.

\section{Minimality via MDL Scoring}
To prove that the generated model $N$ is not an overfitted "flower net," DTEAM embeds a Minimum Description Length (MDL) score directly into the artifact.
\begin{equation}
\Phi(N) = |T| + \left( |A| \cdot \log_2(|T|) \right)
\end{equation}
Where $|T|$ is the number of transitions and $|A|$ is the number of arcs. By demonstrating that $\Phi(N)$ is minimized relative to the trace footprint, DTEAM mathematically proves that its 100\% classification accuracy is achieved via true structural generalization, not memorization.
